General notes
\begin{itemize}
    \item Constantly think about that you are not writing this for people from all disciplines, but specifically for a psychological audience! This audience has a specific frame-of-reference and background knowledge (e.g., SEM, MLM).
    \item Make sure that each section only contains what it should: do not provide discussion of results in results, but save for discussion; do not provide background information in the results, etc. In line with this, changing the research report, use level 2/3 headers after the introduction to introduce background information (e.g., assumptions and criterions).
    \item At the end: check references and rules citation/formatting style (e.g., headers, figure captions, figure notes).
    \item make sure that figures/tables use color coding and aspects only when contributing to the explanation. 
\end{itemize}

Note that GEE is a very general approach framework (see \textcite{zeger_models_1988}), but that we employ the word according to its common use case, where we assume and impose a certain working correlation structure.

it is the within effect, which is commonly of interest and relevant for causal interpretation.

"In the presence of contextual effects, \textcite{Ludtke2008} shows that the observed mean centering method yields biased results, while the latent mean centering method is unbiased" \parencite{asparouhov_latent_2019}

"Intensive longitudinal categorical outcomes present the same
analytic issues that occur for their equivalent continuous outcomes.
For example, if the goal is to investigate causal processes at the
within-subject level, it is important to distinguish between-subjects
from within-subjects variation in putative causal X’s; likewise, it is
important to take account of time T as a potential third variable.
Categorical outcomes, however, also present unique problems. As
we will see, they require the use of nonlinear models where the
outcome is not modeled directly, but rather indirectly, through its
probability." \parencite[]{bolger_modeling_2013}

"for binary outcomes, the partition of variance between different levels does not have the intuitive interpretation of the linear model" \parencite{merlo_brief_2006}. 

\section{empirical example}

\subsubsection{\textbf{1. Interpersonal Conflict and Negative Affect}}

\begin{itemize}
    \item \textbf{Predictor}: Interpersonal conflict
    \begin{itemize}
        \item \textit{Dichotomous}: Conflict occurred that day (yes/no)
        \item \textit{Continuous}: Intensity of conflict (e.g., self-rated, 0–10 scale)
    \end{itemize}
    \item \textbf{Outcome}: Negative affect
    \begin{itemize}
        \item \textit{Dichotomous}: Clinically meaningful distress present (e.g., ≥ threshold item on emotion checklist)
        \item \textit{Continuous}: Composite score of negative affect (e.g., PANAS-NA)
    \end{itemize}
    \item \textbf{Plausibility}:
    \begin{itemize}
        \item \textbf{Within-person}: Daily conflicts predict increased distress the same day (Bolger et al., 1989)
        \item \textbf{Between-person}: Individuals with frequent conflict show higher trait negative affect (Cichy et al., 2009)
    \end{itemize}
\end{itemize}

\subsubsection{\textbf{2. Social Rejection and Self-Esteem}}

\begin{itemize}
    \item \textbf{Predictor}: Experience of rejection
    \begin{itemize}
        \item \textit{Dichotomous}: Rejection occurred (e.g., being excluded, criticized)
        \item \textit{Continuous}: Perceived intensity of rejection
    \end{itemize}
    \item \textbf{Outcome}: Self-esteem
    \begin{itemize}
        \item \textit{Dichotomous}: Low self-esteem moment present (e.g., based on dichotomous item)
        \item \textit{Continuous}: State self-esteem score
    \end{itemize}
    \item \textbf{Plausibility}:
    \begin{itemize}
        \item \textbf{Within-person}: Momentary rejection reduces state self-esteem (Leary et al., 1995)
        \item \textbf{Between-person}: More rejection experiences → lower global self-esteem (Ayduk et al., 2001)
    \end{itemize}
\end{itemize}

\subsubsection{\textbf{3. Prosocial Behavior and Positive Affect}}

\begin{itemize}
    \item \textbf{Predictor}: Engaged in a prosocial act
    \begin{itemize}
        \item \textit{Dichotomous}: Yes/no (e.g., helped someone, gave a compliment)
        \item \textit{Continuous}: Time or effort spent helping
    \end{itemize}
    \item \textbf{Outcome}: Positive affect
    \begin{itemize}
        \item \textit{Dichotomous}: High-mood episode occurred (e.g., endorsement of strong positive emotion)
        \item \textit{Continuous}: Positive affect rating
    \end{itemize}
    \item \textbf{Plausibility}:
    \begin{itemize}
        \item \textbf{Within-person}: Helping behavior boosts mood (Layous et al., 2012)
        \item \textbf{Between-person}: More prosocial individuals report greater well-being (Aknin et al., 2013)
    \end{itemize}
\end{itemize}

\subsubsection{ 4. Intimacy and Conflict}

Chapter 5:

"The example dataset in this chapter is, as in the previous one, simulated
to correspond to datasets from actual intensive longitudinal
studies with which we have been involved (Bolger, DeLongis, Kessler,
& Schilling, 1989; Bolger, DeLongis, Kessler, & Wethington,
1989, 1990; Bolger, Stadler, Paprocki, & DeLongis, 2009; Kennedy,
Bolger, & Shrout, 2002). We simulated a dataset to represent 28 consecutive
daily reports of daily relationship conflicts and daily relational
intimacy from 66 women in a cohabiting, intimate relationship. \ldots 
As in the previous dataset, we simulated scores
on the Reis and Shaver Intimacy Scale (intimacy) to lie on a 0 to 10
interval, such that 0 was the lowest possible score and 10 was the
highest possible score. The variable conflict was simulated such that 0
represented no conflict today, and 1 represented having at least one
conflict today. Finally, we simulated responses to a single-item, cross-sectional
measure of how globally satisfied respondents felt about their intimate
relationship. We simulated responses as follows (with n’s in
parentheses): 5 = best I could ever imagine a relationship being
(n = 10), 4 = extremely satisfied (n = 18), 3 moderately satisfied
(n = 12), 2 = a bit dissatisfied (n = 19), 1 = very dissatisfied (n = 7).
Note that for ease of exposition, relationship quality (relqual) is
treated as a dichotomy of low- and high-quality relationships. A
dichotomous (0, 1) version used in the analyses was created by
combining codes 3, 4, and 5 into a high-relationship quality group
(n = 40, 61%) and codes 1 and 2 into a low-relationship quality
group (n = 26, 39%)" \parencite[]{bolger_intensive_2013}

Chapter 6:

"The dataset is from a daily diary study of couples during what was
for them a typical 4-week period. Each day both partners in each
couple provided reports in the morning, within an hour of rising,
and in the evening, within an hour of going to bed. In the morning,
each partner reported his or her momentary (i.e., “right now”) emotional
states as well as sleep quality and duration. In the evening,
each partner again reported on momentary states and on events that
had occurred during the day. Analyses based on this dataset have
already appeared in several publications (see Gleason, Iida, Bolger,
& Shrout, 2003; Gleason, Iida, Shrout, & Bolger, 2008; Kennedy et
al., 2002). The current analyses are based on a subset of 61 heterosexual
couples from the original dataset.
Here we examine evening reports of conflict from the point of
view of the male partners to see if the anger/irritability of the female
partner at the beginning of a given day increased the risk of such
conflicts later in the day. As is always the case with intensive longitudinal
data, we allow for between-subjects differences (i.e., random
effects), which in this case are between-couples differences, in the
occurrence of conflict and in the anger–conflict relationship"

\subsubsection{\subsubsection{5. Mindfulness and sleep quality}}

\paragraph{\textbf{Predictor: Mindfulness}}

\begin{itemize}
    \item \textbf{Dichotomous (event-based):} Practiced mindfulness today (yes/no)
E.g., engaged in formal mindfulness (meditation, mindful breathing), or reported being mindful (≥ threshold on a dichotomous item)

    \item \textbf{Continuous:} Momentary or daily mindfulness level
E.g., state mindfulness scale (Brown \& Ryan, 2003), such as the Toronto Mindfulness Scale or MAAS daily version

\end{itemize}

\paragraph{\textbf{Outcome: Sleep Quality}}

\begin{itemize}
    \item \textbf{Dichotomous:} Had poor sleep (yes/no)
Based on sleep diary item or threshold on subjective sleep quality (e.g., “slept poorly” vs. not)

    \item \textbf{Continuous:} Subjective sleep quality rating (0–10), sleep onset latency, or sleep efficiency from actigraphy
\end{itemize}

\subsubsection{\textbf{6. Mindfulness and Rumination}}

\begin{itemize}
    \item \textbf{Predictor}: Mindfulness
    \item \textbf{Outcome}: Rumination (or related cognitive interference)
\end{itemize}

\textbf{Within-person}:

When you are more mindful than usual, you ruminate less.

\paragraph{\textbf{Between-person}:}

Individuals who regularly practice mindfulness might have a history of \textbf{higher baseline rumination} (e.g., using mindfulness to cope with it), so their \textbf{mean level of rumination is higher}, despite within-person benefits.

→ \textbf{Result}: Negative within-person, positive between-person → negative contextual effect

\begin{itemize}
    \item \textbf{Theoretical basis}: Selection effect—those prone to rumination seek out and report more mindfulness use.
    \item \textbf{Empirical support}: Cross-sectional studies show mixed associations depending on whether mindfulness is dispositional or practice-based (Barnhofer et al., 2011).
\end{itemize}
 